\section{Conclusion}
\label{sec:conclusion}

% We have investigated whether short-pulse sequences can participate directly in
% superconducting-qubit frequency calibration.  In the example studied here, a
% two-branch sequence provides a local detuning discriminator without a
% free-evolution scan.  A full third-order response kernel improves its local
% inverse, while a drive reference that tracks the predicted qubit frequency
% prevents the feedback trajectory from exhausting that local validity.  In the
% simulated 20-MHz tuning task, the protocol reduces the instrumented solver cost
% by a factor of 4.1 relative to double-sweep Ramsey and reaches a smaller
% analytically checked residual. A coarse short-pulse stage followed by Ramsey
% refinement does not provide the same advantage.

% The result demonstrates a qualified route for bringing short-pulse dynamics
% into calibration: identify a parameter-sensitive pulse response, calibrate its
% local inverse, and preserve that inverse's validity as the device moves.  It
% does not establish that short sequences replace spectroscopy or Ramsey for
% global acquisition.  Experimental validation with finite shots, readout
% error, drift, and control imperfections will determine how broadly this role
% extends and whether it reduces acquisition time on hardware.
% The extended \textsc{Acquire--Track--Verify--Lock} architecture provides a
% testable route to that validation, but the present results establish only the
% deterministic local-tracking case, not hardware verification or long-term
% locking performance.
